% =====================================================================
% Appendix B --- Dataset Details and Additional Results
% =====================================================================

This appendix records the configuration, training, and filter parameters needed
to reproduce the experiments of Chapters~\ref{cha:monocular} and~\ref{cha:airsim}.

\section{AirSim Dataset Configuration}
\label{ap:airsim_cfg}

\paragraph*{Environments and sensors.} The synthetic data of
Chapter~\ref{cha:airsim} were generated in three Unreal/AirSim environments,
\textsc{AirSimNH} (low-rise residential), \textsc{City} (dense downtown), and
\textsc{Coastline} (vegetated coastal). The ego platform carries an RGB camera, a
planar-depth camera, and a $64$-beam LiDAR, rigidly co-located ($0.35$~m forward,
$0.5$~m below the body origin) and pitched $-15^\circ$ so that they share one
optical frame. The camera is an ideal pinhole at $1280\times720$ with a
$90^\circ$ horizontal field of view ($f_x=f_y=640$~px, principal point
$(640,360)$~px, no lens distortion). Depth returns are kept for
$0.1\le z<250$~m, yielding $10^{5}$--$3\times10^{5}$ points per frame; the depth
camera, not the LiDAR, is the primary cloud, because at altitude the targets are
seen against the sky where the LiDAR returns nothing.

\paragraph*{Targets and 3D annotation.} Each target carries a visual quadrotor
mesh (\emph{VisQuad}) enlarged $\times4$ so as to be visible and densely sampled
at $20$--$100$~m; its axis-aligned full extents are $3.01\times3.93\times2.79$~m,
which after a $1.38$ margin (for the omitted spinning rotors) and halving give the
annotation half-extents $\bm{e}=[2.08,2.71,1.93]^{\top}$~m. The ground-truth 3D box
is built geometrically from the true global pose (not from the simulator's
detection API, which carries a systematic vertical offset). Per-frame acquisition
is performed with the scene frozen (\texttt{simPause}) and targets/ego
re-teleported to their planned poses, to keep image and reported poses consistent.

\paragraph*{Training and evaluation subsets.} The 2D detector is trained on
$\approx2{,}500$ pixel-exact synthetic-composition frames ($196$ drone sprites
spanning $12$~yaw~$\times$~$3$~pitch~$\times$~$3$~roll orientations at five ranges,
$1$--$5$ drones per frame). The point-cloud networks are trained on the collected
multimodal set ($3$ scenes $\times\,1{,}000$ frames). Evaluation uses a disjoint
campaign of $60$ video sequences ($20$ per environment, $75$--$110$ frames at
${\approx}5$~Hz) that randomizes the range regime (near $18$--$32$~m, mid
$32$--$50$~m, far $50$--$78$~m), time of day ($07$h--$18$h), weather (clear, light
fog, light rain), and number of targets ($2$--$3$), with a \emph{moving} observer.

\section{Network Training Hyperparameters}
\label{ap:train_hyper}

Table~\ref{tab:hyper_yolo} lists the 2D-detector training settings and
Table~\ref{tab:hyper_3d} the point-cloud--network settings, for reproducibility.

\begin{table}[H]
  \centering
  \footnotesize
  \caption{YOLOv11-small training (Ultralytics defaults unless noted): SGD with
  Nesterov momentum, weight decay, and a cosine learning-rate schedule.}
  \label{tab:hyper_yolo}
  \resizebox{\linewidth}{!}{%
  \begin{tabular}{@{}llll@{}}
    \toprule
    Setting & Synthetic drone (Ch.~\ref{cha:airsim}) & Real-urban drone (ref.) & Person (Ch.~\ref{cha:monocular}) \\
    \midrule
    Backbone        & YOLOv11s ($\approx 10$~M) & YOLOv11s & YOLOv11s \\
    Input size      & $1280$~px & $640$~px & $1280$~px \\
    Epochs          & $80$ & $100$ & $50$ \\
    Batch           & $16$ & $16$ & auto \\
    Patience        & default & $30$ & $100$ \\
    Pretrained      & COCO & COCO & COCO \\
    Augmentation    & mosaic, HSV, scale & mosaic, HSV, scale & Ultralytics default \\
    Result ($\mathrm{mAP}_{50}$) & $0.994$ & $0.731$ & $0.754$ \\
    \bottomrule
  \end{tabular}}
\end{table}

The Chapter~\ref{cha:monocular} person detector (rightmost column) was trained
through the Ultralytics HUB on the dataset \texttt{visdrone\_sard\_heridal}, the
union of the VisDrone, SARD, and HERIDAL aerial person-detection sets:
$14{,}783$ images of a single \texttt{person} class, split $74.1\%/12.4\%/13.5\%$
into train/validation/test ($10{,}961$, $1{,}829$, and $1{,}993$ images). Beyond
the $\mathrm{mAP}_{50}=0.754$ in the table, it reaches
$\mathrm{mAP}_{50\text{--}95}=0.389$, with precision $0.795$ and recall $0.669$ on
its validation split.

\begin{table}[H]
  \centering
  \footnotesize
  \caption{Point-cloud detector training. All use Adam ($10^{-3}$, weight decay
  $10^{-4}$) with a cosine schedule; frustum nets sample $N{=}4096$ points from
  the $25\%$-enlarged frustum and normalize by centroid subtraction and a scale.}
  \label{tab:hyper_3d}
  \begin{tabular}{@{}p{2.5cm} p{6.3cm} p{5.3cm}@{}}
    \toprule
    Network & Architecture (key) & Loss / schedule \\
    \midrule
    Frustum PointNet\texttt{++}
      & SA $[64,64,128]$ / $[128,128,256]$ / $[256,512,1024]$, FC $[512,256,128]$
      & Focal ($\alpha{=}0.5,\gamma{=}2$) $+\,5\,\mathrm{SL}_1(\text{ctr})+\mathrm{SL}_1(\text{size})$; $30$ ep., batch $8$ \\[2pt]
    T-Net segmentation
      & $+$ per-point fg/bg head ($1088{\to}256{\to}128{\to}1$)
      & $+\,\mathrm{BCE}(\text{seg})$; $20$ ep., batch $32$ \\[2pt]
    Frustum ConvNet
      & $16$ depth slices, per-slice mini-PointNet, $1$D conv
      & same heads; $30$ ep., batch $8$ \\[2pt]
    PointPillars (PFN)
      & BEV $1.25$~m grid, per-pillar mini-PointNet $[64,64]$, 2D CNN
      & center heatmap (focal) $+$ box reg.; $15$ ep., batch $8$ \\[2pt]
    Dense 3D voxel
      & $2.0$~m grid, Conv3D $1{\to}16{\to}32{\to}32{\to}64$
      & center heatmap $+$ box reg.; $15$ ep., batch $4$ \\
    \bottomrule
  \end{tabular}
\end{table}

\section{Filter Parameter Settings}
\label{ap:filter_params}

\paragraph*{3D range--bearing EKF (Chapters~\ref{cha:tracking} and~\ref{cha:airsim}).}
Acceleration intensity $\sigma_a=0.5~\mathrm{m/s^2}$; bearing noise
$\sigma_\alpha=\sigma_\beta=5~\mathrm{mrad}$; range noise $\sigma_\rho=0.05\,\rho$
(floored at $0.5$~m); size noise $\sigma_r=0.5$~m; initial position/velocity
standard deviations $5$~m / $10~\mathrm{m/s}$; gating at the $\chi^2$ $99\%$ level
($13.28$ for four degrees of freedom). The SORT-style association uses
non-maximum-suppression IoU $0.4$, association-gate IoU $0.3$, minimum hits $3$,
coasting horizon $4$, and maximum age $8$ frames (reproduced from
Table~\ref{tab:tracker}).

\paragraph*{Monocular ground-plane EKF (Chapter~\ref{cha:monocular}).}
Planar constant-velocity model with acceleration intensity $\sigma_a$; pixel
measurement noise $\bm{R}_{\mathrm{cam}}=\sigma_{\mathrm{cam}}^2\bm{I}_2$; the
$6$-DoF pose covariance is modeled as isotropic,
$\bm{\Sigma}=\sigma_{\mathrm{pose}}^2\bm{I}_6$, with $\sigma_{\mathrm{pose}}$ tuned
by the NIS consistency analysis of Section~\ref{sec:mono_consistency} to
$\sigma_{\mathrm{pose}}\approx10^{-3}$, the value at which the mean NIS crosses the
consistent value $\bar\eta=2$. The first detection is initialized by ray--plane
intersection with a large initial covariance so that early measurements dominate.

\section{Code and Data Availability}
\label{ap:code}

The synthetic-dataset generation and 3D-detection pipeline of
Chapter~\ref{cha:airsim} are released as open source at
\url{https://github.com/ericyoshida/airsim-uav-3d-detection}. The monocular
tracker of Chapter~\ref{cha:monocular} is available from the author on request;
its public release, together with any TU~Berlin flight data, is subject to the
consent of the collaborating team. In both cases the large artifacts (datasets,
model weights, recordings) are kept separate from the code.
